\documentclass[11pt,letterpaper]{article}
\usepackage[margin=1in]{geometry}
\usepackage[T1]{fontenc}
\usepackage{lmodern,microtype}
\usepackage{amsmath,amssymb,amsthm,mathtools}
\usepackage{booktabs,longtable,array}
\usepackage{enumitem}
\usepackage{xcolor}
\usepackage{xurl}
\usepackage{hyperref}
\usepackage{fancyhdr}
\usepackage{listings}
\definecolor{ink}{HTML}{183A51}
\definecolor{linkblue}{HTML}{246188}
\hypersetup{colorlinks=true,linkcolor=linkblue,urlcolor=linkblue,citecolor=linkblue,
 pdfauthor={AI-assisted research prepared for Vladimir Reshetnikov},
 pdftitle={Factorial fixed points and smaller Diophantine consistency equations}}
\pagestyle{fancy}
\fancyhf{}
\fancyhead[L]{\small\textcolor{ink}{Diophantine consistency equations}}
\fancyhead[R]{\small October 2026}
\fancyfoot[C]{\thepage}
\setlength{\headheight}{14pt}
\setlength{\emergencystretch}{2em}
\setlength{\parskip}{3pt}
\setlist{itemsep=3pt,topsep=5pt}
\lstset{basicstyle=\small\ttfamily,breaklines=true,columns=fullflexible,
 keepspaces=true,frame=single,rulecolor=\color{black!18},xleftmargin=3pt,xrightmargin=3pt}
\newtheorem{theorem}{Theorem}[section]
\newtheorem{lemma}[theorem]{Lemma}
\newtheorem{proposition}[theorem]{Proposition}
\newtheorem{corollary}[theorem]{Corollary}
\theoremstyle{definition}
\newtheorem{definition}[theorem]{Definition}
\theoremstyle{remark}
\newtheorem{remark}[theorem]{Remark}
\newcommand{\N}{\mathbb N}
\newcommand{\Z}{\mathbb Z}
\newcommand{\Q}{\mathbb Q}
\newcommand{\Con}{\operatorname{Con}}
\newcommand{\ZFC}{\mathrm{ZFC}}
\newcommand{\ZF}{\mathrm{ZF}}
\newcommand{\htp}{\operatorname{ht}}
\newcommand{\ceil}[1]{\left\lceil#1\right\rceil}
\newcommand{\file}[1]{\nolinkurl{#1}}
\newcommand{\proofpred}{\operatorname{Proof}_{\ZFC}}
% Final metrics from exact coefficient files.
\newcommand{\NatVars}{468}
\newcommand{\NatRows}{364}
\newcommand{\NatTerms}{14\,842}
\newcommand{\NatHeight}{83\,850\,989}
\newcommand{\IntVars}{956}
\newcommand{\IntRows}{510}
\newcommand{\IntTerms}{10\,988}
\newcommand{\IntHeight}{2\,445\,790\,974\,704}
\newcommand{\SmallVars}{989}
\newcommand{\SmallRows}{543}
\newcommand{\SmallTerms}{7\,705}
\newcommand{\SmallHeight}{63\,910\,295}
\newcommand{\ParentVars}{781}
\newcommand{\ParentRows}{677}
\newcommand{\SingleGuards}{163}
\newcommand{\PairGuards}{8}
\newcommand{\GuardedVars}{1294}
\newcommand{\GuardedRows}{848}
\newcommand{\ProjectSteps}{338}
\newcommand{\SeedRules}{610}
\newcommand{\RepoPin}{8f8fac484b976ba0ab75fd173282e838ca5a3b85}


\begin{document}
\hypersetup{pageanchor=false}
\begin{titlepage}
\vspace*{0.55in}
{\large\sffamily\color{ink} A research submission for the ProveIt repository\par}
\vspace{0.35in}
{\Huge\bfseries Factorial fixed points\par}
\vspace{0.10in}
{\huge\bfseries and smaller Diophantine\par consistency equations\par}
\vspace{0.35in}
{\Large Explicit quartics, complete coefficient files,\par and shorter small-variable circuits\par}
\vspace{0.45in}
{\large Prepared for Vladimir Reshetnikov\par}
{\large AI-assisted mathematical exploration and exact algebra verification\par}
{\large 24 October 2026\par}
\vfill
\begin{abstract}
We give several completely materialized polynomial equations whose absence
of solutions is equivalent in ZFC to $\Con(\ZFC)$, by composing new
arithmetic equivalences with the repository's identified proof-calculus
calibration. The principal new lemma characterizes $n!$ as the unique
fixed point
\[
 F=\ceil{\frac{(F^2+1)^n}{\binom{F^2+n}{n}}},\qquad n,F\in\N.
\]
It permits a proposed factorial value to supply its own precision bound.
One square replaces a separate variable-exponent radius certificate.
Together with a smaller local rule certificate and exact sign inference,
this gives a quartic in \IntVars\ unrestricted integer variables with
\IntRows\ quadratic residuals, \IntTerms\ monomials, and explicitly listed
coefficients. A companion in \SmallVars\ integers has \SmallTerms\
monomials and coefficient height below $2^{26}$.
Natural-domain variants use 467, 468, or 469 variables. Independent
improvements give a 456-operation nine-positive-variable equation of
degree $1\,376\,378$ and shorter exact radical norms.
The archive contains all literal equations, proofs, generators,
independent algebra checkers, source pins, and verification receipts.
No global minimum, worldwide priority, proof of consistency, or new
end-to-end proof-assistant certificate is asserted.
\end{abstract}
\vspace{0.25in}
{\small Comparison snapshot: \texttt{\RepoPin}.\par
All variable counts refer to the domain stated next to the equation.}
\end{titlepage}
\hypersetup{pageanchor=true}
\begingroup
\small
\setlength{\parskip}{0pt}
\tableofcontents
\endgroup
\newpage

\section{What has been constructed}
\label{sec:results}

There are several reasonable meanings of ``simplest'': few unknowns,
low degree, a short arithmetic circuit, few expanded terms, small
coefficients, or a short semantic proof. The delivered equations improve
several of these measures, but they do not order all the measures at once.
The main construction follows the direct C222 proof calculus. It uses
neither a simulation of the K215 Turing machine nor Jones' universal
polynomial. The separate small-variable improvements retain their own
parent source encodings and are identified as such.

Throughout, $\N=\{0,1,2,\ldots\}$, $\Z_{>0}=\{1,2,\ldots\}$,
and an unqualified integer variable ranges over all of $\Z$.
The height of an expanded integer polynomial is the maximum absolute
value of its collected coefficients. A quadratic residual system
$r_1=\cdots=r_s=0$ specifies the \emph{single} polynomial
\begin{equation}\label{eq:sos-definition}
                         P=\sum_{j=1}^{s}r_j^2.
\end{equation}
The residuals are polynomial abbreviations, not additional unknowns.
Every listed coordinate, including every retained arithmetic gate, is
included in the variable count.

\begin{theorem}[Main explicit integer equations]\label{thm:main}
Let $P_{956}$ be the polynomial specified by
\file{integer/equations/integer956/polynomial.json}; equivalently it is
the sum of the squares of the 510 literal residuals in the adjacent
\file{equation.json}. It belongs to $\Z[X_0,\ldots,X_{955}]$, has exact
degree four, 10,988 nonzero monomials, and height
$2\,445\,790\,974\,704$. The ordinary mathematical proof in this
submission, with its identified inherited calibration, proves in ZFC
\begin{equation}\label{eq:main-con}
 \Con(\ZFC)\quad\longleftrightarrow\quad
 \forall (x_0,\ldots,x_{955})\in\Z^{956},\quad P_{956}(x)\ne0.
\end{equation}
The companion $P_{989}$ in
\file{integer/equations/integer989/polynomial.json} satisfies the same
biconditional with 989 integer unknowns. It has 543 quadratic residuals,
7,705 monomials, exact degree four, and height $63\,910\,295<2^{26}$.
\end{theorem}

The proof is constructive at the level of finite encodings: a zero
recovers a finite ZFC refutation, and a refutation supplies a zero.
Thus the more informative form of the result is
\[
 \exists p\in\N\ \proofpred(p,\bot)
 \quad\longleftrightarrow\quad
 \exists x\in\Z^{956}\ P_{956}(x)=0.
\]
Negating this equivalence gives \eqref{eq:main-con}.
No choice of coefficients depends on assuming consistency or
inconsistency. No actual zero or proof of the absence of zeros is supplied.

\subsection{The direct quartic tradeoffs}

Table~\ref{tab:quartics} compares the new direct family with the latest
relevant incoming direct constructions at the pinned snapshot. The
older 2,088-integer conversion is \emph{not} used as the principal integer
baseline: the incoming batched-certificate construction already has
1,690 integer variables. The parent and domain of each comparison matter.

\begin{table}[htbp]
\centering\small
\begin{tabular}{@{}llrrrr@{}}
\toprule
Construction & Domain & Variables & Rows & Monomials & Height\\
\midrule
Synchronized incoming & $\N$ & 522 & 400 & 20,225 & 83,900,628\\
New, output-linear radius & $\N$ & 467 & 363 & 18,479 & 83,880,005\\
New, square radius & $\N$ & 468 & 364 & 14,842 & 83,850,989\\
New, smaller support & $\N$ & 469 & 365 & 12,051 & 83,850,757\\
\addlinespace
Batched incoming & $\Z$ & 1,690 & 842 & 21,381 & 2,649,643,805,758\\
New, fewer variables & $\Z$ & 956 & 510 & 10,988 & 2,445,790,974,704\\
New, smaller coefficients & $\Z$ & 989 & 543 & 7,705 & 63,910,295\\
\bottomrule
\end{tabular}
\caption{Direct literal equations. Every degree is exactly four; ``rows''
counts the quadratic residuals before the final squared sum.}
\label{tab:quartics}
\end{table}

The 467-variable natural equation uses the radius $H^2F_H$; the 468-
and 469-variable equations use $F_H^2$. The latter has a smaller raw
parent and fewer expanded terms, although the deterministic elimination
order leaves one more variable. This is a real tradeoff, so both complete
variants are retained.

Relative to the 1,690-variable direct integer baseline, $P_{956}$ removes
734 variables and 10,393 monomials. These are reductions of about
43.4\% and 48.6\%, respectively. It is not a claim that 956 is the
smallest possible integer-variable count: the repository already contains
much smaller variable counts with different degrees and coefficient
representations.

\subsection{Separate improvements at small variable counts}

\begin{table}[htbp]
\centering\small
\begin{tabular}{@{}lrrrr@{}}
\toprule
Domain & Variables & Degree & Parent operations & New operations\\
\midrule
$\Z_{>0}$, original31 & 9 & 1,376,378 & 475 & 456\\
$\Z_{>0}$, compact54 & 9 & 1,520,186 & 433\textsuperscript{a} & 442\\
$\Z_{>0}$ & 11 & 211 & 111 & 110\\
$\Z$ & 13 & 1,408 & 185 & 176\\
$\Z$ & 13 & 1,392 & 186 & 177\\
$\Z$ & 13 & 1,088 & 195 & 186\\
$\Z$ & 13 & 1,072 & 195 & 186\\
\bottomrule
\end{tabular}
\caption{Live binary arithmetic operations. Addition, subtraction,
multiplication, squares, and scalar multiplication are charged; register
reuse is free. \textsuperscript{a}The 433-operation parent has degree
1,628,768, so this row is a cost--degree tradeoff.}
\label{tab:small}
\end{table}

For the first row, the previous degree is $1\,379\,730$: both cost and
degree improve. For the last five rows the entire polynomial is unchanged;
only its evaluator is shorter. In particular, the degree-1,072 result
improves the latest incoming 195-operation degree-1,072 equation, rather
than claiming its parent's degree reduction anew.

These arithmetic operation counts treat a fixed numeral as a leaf.
The fixed leaves in the small-variable families can be very large.
Their complete bytes and reconstruction bindings are included, but the
cost measure is not a bit-operation or printed-digit complexity measure.
The direct quartics instead provide ordinary, fully collected, modest-sized
integer coefficients.

\section{The factorial fixed point}
\label{sec:factorial}

The usual ceiling representation of a factorial chooses an independent
radius large enough to make a rational approximation lie within one of
$n!$. The new observation is that the unknown value can help select the
radius, provided the resulting fixed point is proved unique. This removes
an arithmetic graph rather than merely abbreviating its notation.

For $n,R\in\N$, define
\[
 f_n(R)=\frac{(R+1)^n}{\binom{R+n}{n}}
       =n!\prod_{j=1}^{n}\frac{R+1}{R+j}.
\]
Empty products are one. The denominator is positive.

\begin{lemma}[Approximation and monotonicity]\label{lem:defect}
For $n\ge2$ and real $R>0$, $f_n(R)$ is positive, nondecreasing in $R$,
at most $n!$, and satisfies
\begin{equation}\label{eq:defect}
           0\le n!-f_n(R)<\frac{n!\,n(n-1)}{2R}.
\end{equation}
The formula for real $R$ here means the finite product, so no generalized
binomial coefficient is needed.
\end{lemma}

\begin{proof}
Each factor is positive and at most one. For $j\ge1$, the function
$(R+1)/(R+j)$ is nondecreasing, by cross multiplication (or differentiation).
Put $u_j=(j-1)/(R+j)$. Then $0\le u_j<1$, and induction on the finite
product gives
\[
 1-\prod_{j=1}^n(1-u_j)\le\sum_{j=1}^n u_j
       <\frac1R\sum_{j=1}^n(j-1)=\frac{n(n-1)}{2R}.
\]
The strict inequality holds because $n\ge2$ gives a positive summand
with denominator strictly greater than $R$. Multiply by $n!$.
\end{proof}

\begin{theorem}[Square-radius fixed point]\label{thm:fixedpoint}
For all $n,F\in\N$,
\begin{equation}\label{eq:fixedpoint}
 F=\ceil{\frac{(F^2+1)^n}{\binom{F^2+n}{n}}}
                   \quad\longleftrightarrow\quad F=n!.
\end{equation}
\end{theorem}

\begin{proof}
Let $K=n!$ and suppose first that $F=\ceil{f_n(F^2)}$.
The quotient is positive and at most $K$, so $1\le F\le K$.
For $n\ge4$, monotonicity and $F^2\ge1$ give
\[
 f_n(F^2)\ge f_n(1)=\frac{2^n}{n+1}>n-1.
\]
The last inequality is $2^n>n^2-1$. It holds at $n=4$; its induction
step follows from
$2(n^2-1)-((n+1)^2-1)=n^2-2n-2>0$ for $n\ge4$.
Consequently $F\ge n$.

Apply \eqref{eq:defect} with $R=F^2$. First,
\[
 K-f_n(F^2)<\frac{K n(n-1)}{2F^2}<\frac K2,
\]
so $F\ge f_n(F^2)>K/2$. Applying the same estimate a second time,
\[
 K-f_n(F^2)
 <\frac{K n(n-1)}{2F^2}
 <\frac{2n(n-1)}K
 =\frac{2}{(n-2)!}\le1.
\]
Thus $K-1<f_n(F^2)\le K$, and its ceiling is $K$.

At $n=0,1$ the quotient is identically one. At $n=2$ and $F\ge1$,
\[
 f_2(F^2)=\frac{2(F^2+1)}{F^2+2}\in(1,2),
\]
so $F=2$. At $n=3$, the already established bound $1\le F\le6$
leaves six exact rational values:
\[
\begin{array}{c|rrrrrr}
 F&1&2&3&4&5&6\\\hline
 \ceil{f_3(F^2)}&2&4&5&6&6&6.
\end{array}
\]
Only $F=6$ is a fixed point. This proves soundness in every case.

Conversely, set $F=K$. For $n\ge2$, \eqref{eq:defect} gives
\[
 0\le K-f_n(K^2)<\frac{n(n-1)}{2K}
                  =\frac{1}{2(n-2)!}\le\frac12.
\]
Hence its ceiling is $K$. The cases $n=0,1$ are immediate.
\end{proof}

This is an unrestricted elementary proof, not an experimental fixed-point
search. The bounded calculations in the audit scripts check transcription
and small cases; they are not the reason the theorem holds for arbitrary
$n$.

\subsection{A common radius for several factorials}

\begin{corollary}[Shared radius]\label{cor:common-radius}
Suppose $H\ge1$, the first ceiling certificate uses the argument $H$
and radius $R=F_H^2$, and every later argument $n$ satisfies $0\le n\le H$.
Then the first certificate forces $F_H=H!$, and the same radius correctly
represents every later factorial by $F_n=\ceil{f_n(R)}$.
\end{corollary}

\begin{proof}
Theorem~\ref{thm:fixedpoint} handles the first certificate. For
$2\le n\le H$, with $R=(H!)^2$,
\[
 0\le n!-f_n(R)
 <\frac{n!\,n(n-1)}{2(H!)^2}
 \le\frac{H(H-1)}{2H!}
 =\frac{1}{2(H-2)!}\le\frac12.
\]
The cases $n=0,1$ have quotient one, and if $H=1$ these are the only
cases. Completeness uses the actual factorial values and their remainders.
\end{proof}

The concrete compiler requests $H!$ first and $t!$ second, with
$H\ge t\ge1$. It uses one square $F_H^2$ for both.
That same square already occurs in the base-modulus equation for the
power request $F_H^t$, so it is shared in the final arithmetic circuit.

\subsection{The exact polynomial ceiling graph}

Let $N,C$ be positive naturals and let $F,r,s$ be naturals. Then
\begin{equation}\label{eq:ceiling-graph}
       N+C=CF+r+1,\qquad C=r+s+1
\end{equation}
is equivalent, existentially in $r,s$, to $F=\ceil{N/C}$.
Indeed the second row gives $0\le r<C$, and the first gives
$(F-1)C<N\le FC$. Conversely these strict and weak inequalities provide
$r=N+C-CF-1\in[0,C-1]$ and $s=C-r-1$.

In the factorial graph, $C=\binom{R+n}{n}$ and $N=(R+1)^n$.
Their notation describes the semantics of separate finite polynomial
graphs. No ceiling, factorial, binomial symbol, or variable exponent
remains in \file{equation.json} or \file{polynomial.json}.

\subsection{The retained output-linear alternative}

An especially short fixed-point proof uses $R=n^2F$:
\begin{equation}\label{eq:linear-radius}
 F=\ceil{f_n(n^2F)}\quad\longleftrightarrow\quad F=n!.
\end{equation}
For $n\ge2$, a proposed fixed point has $1\le F\le K=n!$, and
\[
 K-f_n(n^2F)<\frac{K(n-1)}{2nF}<\frac{K}{2F}.
\]
Initially $F\ge1$ forces $f_n(n^2F)>K/2$, hence $F>K/2$.
The same inequality then gives defect less than one, proving soundness.
For $F=K$ the defect is less than $1/2$, proving completeness.
The cases $n=0,1$ are immediate. With a dominant argument $H$, the
shared radius $H^2F_H=H^2H!$ gives, for $2\le n\le H$, a defect
strictly below $(H-1)/(2H)<1/2$.

This alternative produces a 783-variable parent and the 467-variable
natural contraction. The square-radius version has a 781-variable parent
and better term counts. Both are compile-time choices of a single literal
polynomial, not choices made by solutions of the equation.

\section{A smaller direct certificate for the proof rules}
\label{sec:source}

The source is the inherited numerical calculus C222, with truth label
$0$ and target falsity label $1$. Its numerical constructors include
\begin{align*}
 I(a,b)&=2(a+b)(a+b+1)+4b,& A(a,b)&=I(a,b)+2,\\
 M(a,b)&=I(a,b)+1,& E(a,b)&=I(a,b)+3,\\
 N(a)&=I(a,1),&X(a,b)&=N(A(a,N(b))).
\end{align*}
These are literal polynomials. Their complete axiom instances occur in
\file{direct/code/hybrid_source.py}, and the constructor prefix is
checked by exact syntax comparison with the inherited source. The
calculus contains truth, eight unconditional axiom schemata, one guarded
bundle schema, modus ponens, and generalization.

\subsection{Eliminate premise labels inside each rule branch}

Earlier local certificates carried independent premise labels $p,q$
even though each selected branch determines them. The new bounded row
uses just nine coordinates
\[
                      (v,a,b,c,d,S,S',u,u').
\]
Here $v$ is the conclusion, $a,b,c,d$ are schema parameters, $S,S'$ are
successive accumulator values, and $u,u'$ are beta quotients.
Let $h_1,\ldots,h_8$ be the eight inherited unconditional axiom
polynomials and $h_B$ the bundle polynomial. Put
\begin{align*}
 G&=((c-a)(c-b))^2-(d+1)^2,\\
 K&=v\prod_{j=1}^{8}(v-h_j),&
 B&=(v-h_B)^2+G^2,\\
 w(z)&=1+z\mathcal M,&\mathcal M&=M_0+1.
\end{align*}
Define the rule residuals and nonnegative branch tests by
\begin{align*}
 L_A&=S'w(v)-S,& Q_A&=K^2B+L_A^2,\\
 L_M&=S'w(v)-S w(a)w(I(a,v)),\\
 L_G&=S'w(v)-S w(b),&Q_G&=(v-A(a,b))^2+L_G^2.
\end{align*}
The local polynomial at index $i$ is
\begin{align}\label{eq:local-source}
 P_{\mathrm{loc}}={}&Q_A L_M^2 Q_G\notag\\
 &+[x-S-(1+(i+1)y)u]^2\notag\\
 &+[x-S'-(1+(i+2)y)u']^2.
\end{align}
Every displayed branch factor is nonnegative, including on real inputs.

If $Q_A=0$, the conclusion is truth or one of the valid axiom instances;
restore $p=q=0$. If $L_M=0$, restore $p=a$, $q=I(a,v)$: this is
modus ponens with the correct accumulator equation. If $Q_G=0$, restore
$p=b$, $q=0$, with $v=A(a,b)$: this is generalization. Conversely each
valid old branch has this parameterization. Unused parameters can be
zero. The proof makes no claim that parameters satisfy unselected
branches. Existential quantification distributes over the finite
disjunction represented by the product.

The guard $G=0$ says
$|(c-a)(c-b)|=d+1>0$, exactly the required pair of disequalities.
Its natural witness exists whenever those disequalities hold.
The branch-specific substitution therefore preserves the actual
Horn rule relation, not merely some numerical evaluations of its source.

\subsection{Restore the correct carrier bound}

Bound $v,a,b,c,d$ by $m\ge1$. The restored second premise in modus
ponens need not be bounded by $m$; it is bounded by $I(m,m)$. Hence use
the label bound
\[
 L(m)=I(m,m)+2=8m^2+8m+2=2(2m+1)^2
\]
and impose
\begin{equation}\label{eq:newcarrier}
                       \mathcal M=(L(m)+1)^t,\qquad t\ge1.
\end{equation}
The inherited bounded-factor carrier theorem applies with $L(m)$.
It converts the product identity at this sufficiently large carrier into
the formal label-multiset identity used by the finite Horn-flow theorem.
All active labels are in its stated bound; large evaluations in inactive
axiom branches need not be label-bounded.

The remaining local bounds are $S,S'\le y$ and $u,u'\le x$.
They make the two beta reconstruction rows in \eqref{eq:local-source}
have strict remainder bounds automatically. The endpoints are
\begin{align*}
 x&=M_0+2+(y+1)u_0,&y&=M_0+r_0+2,\\
 x&=1+[1+(t+1)y]u_t.
\end{align*}
Thus the initial accumulator is $\mathcal M+1$ (target label one),
the final accumulator is one, and the beta code links consecutive
accumulators. Their product telescopes to the certified Horn flow.

Given a finite proof, first choose sufficiently large $t,m$, then the
carrier, the finite accumulator path, and beta data $y,x$. These are
ordinary finite objects. No uniform bound on a possible ZFC refutation
is assumed or needed.

\subsection{The explicit majorant}

Split the signed arithmetic circuit for $P_{\mathrm{loc}}$ into
nonnegative-coefficient circuits $E-F$. Addition, subtraction, and
multiplication on pairs use the exact ring identities
\[
 (e,f)+(g,h)=(e+g,f+h),\quad
 (e,f)-(g,h)=(e+h,f+g),
\]
\[
                    (e,f)(g,h)=(eg+fh,eh+fg).
\]
Evaluate at the index bound $t$ and coordinate bounds
$(m,m,m,m,m,y,y,x,x)$. The actual stored DAG has a checked power
majorant exponent $6\,893\,174$. Set
\begin{equation}\label{eq:H}
 H=(2+t+x+y+m+M_0)^{2^{23}}.
\end{equation}
Since $2^{23}=8\,388\,608$ is larger, $H$ bounds both sides, the
index, and all local bounds with the required strict margin.
Equation~\eqref{eq:H} takes exactly 23 successive squarings.

The certificate checks this bound gate by gate. If the base is
$R\ge2$, input atoms are at most $R$, a sum of quantities bounded by
$R^a,R^b$ is at most $R^{\max(a,b)+1}$, and a product is at most
$R^{a+b}$. Constants are bounded by a suitable fixed power of two.
No semantic claim about a refutation is hidden in this numerical
majorant.

\section{Inherited compiler contracts and the consistency endpoint}
\label{sec:inherited-compiler-contracts}

The new projection and factorial argument are composed with the following
specific arithmetic interfaces.  All variables in this section range over
$\mathbb N=\{0,1,\ldots\}$ unless indicated otherwise.  For the unrestricted
integer equation, the sign-recovery certificate must first establish this
natural domain.  The statements below concern every satisfying assignment;
their use does not presume the intended values of intermediate witnesses.

\subsection{Logical calibration and what is imported}

The numerical source is the fixed twelve-branch calculus $C_{222}$, with
truth code $0$ and falsity code $1$.  Its calibration is
\[
 C_{222}(1)
 \ \longleftrightarrow\
 \operatorname{BProv}(\mathrm{ZF},\bot)
 \ \longleftrightarrow\
 \operatorname{BProv}(\mathrm{ZFC},\bot).
\]
The exact named results are
\nolinkurl{KEndpoint.c222_iff_zf_refutation} and
\nolinkurl{ConstructibleChoice.zf_refutation_iff_zfc_refutation}.
Their frozen files, relative to the direct-equation component, are
\nolinkurl{provenance/upstream/calibration/KEndpoint.lean} and
\nolinkurl{provenance/upstream/calibration/ConstructibleChoice.lean}.
The source numerals and rules are in the adjacent
\nolinkurl{Code254.lean} and \nolinkurl{KCode222.lean}.
The repository paths are
\nolinkurl{Computability/BusyBeaver/Lean/BusyBeaver/ZFC/}, at calibration
commit \nolinkurl{08200eb10fde2db408c69dffd840d6920a569b1f}; exact file hashes
are retained in the provenance ledgers.

These are substantial imported mathematical results.  The first combines
the soundness of the numerical calculus with its model-existence
completeness argument; the second removes Choice by the constructible-model
argument and first-order soundness/completeness.  No Turing-machine endpoint
is used.  The appeal is to these theorems and their proofs formalizable in
ZFC, not to an inference that every external Lean theorem is automatically
a theorem of ZFC.  The complete transitive Lean development was not rebuilt
for this contribution, and no new object-language ZFC proof term is claimed.

The remaining passage from finite rule instances to a proof is the
single-conclusion Horn-flow criterion.  Give $(p,q)\to v$ the vector
$e_v-e_p-e_q$, with $e_0=0$.  A proof tree has total flow $e_1$.
Conversely, if a finite nonnegative sum of rule vectors equals $e_1$,
sum coordinates outside the least rule-closed set containing $0$.
Each rule contributes at most zero outside that set; a missing target
would contribute one, a contradiction.  Cycles do not invalidate this
argument.

For the accumulator-to-flow step, all restored labels lie in
$[0,L]$, where $L=8m^2+8m+2$, and $M=(L+1)^t$, $t\ge1$.
Telescoping the accumulator equations with endpoints $S_0=1+M$, $S_t=1$
gives $P(M)=Q(M)$, where
\[
 P(X)=\prod_{i<t}(1+v_iX),\qquad
 Q(X)=(1+X)\prod_{i<t}(1+p_iX)(1+q_iX).
\]
Since $P(M)<M^{t+1}$, nonnegativity first forces $\deg Q\le t$.
Both products then have at most $t$ nonconstant factors, so every
coefficient is strictly below $M$.  The least nonzero coefficient of
$P-Q$ could not vanish on evaluation at $M$; hence $P=Q$.
Unique factorization in $\mathbb Q[X]$ gives precisely the required flow.
This is the bounded-factor argument in
\nolinkurl{provenance/upstream/calibration/direct_flow_large_carrier_README.md}.
Conversely, expand a proof tree and let $S_i$ be the product of the weights
of its open obligations.  A finite CRT beta code gives
$x=S_i+(1+(i+1)y)u_i$, choosing $y$ larger than every state and divisible
by $1,\ldots,t+1$.  In soundness, the supplied bounds
$S_i\le y<1+(i+1)y$ make these actual remainders and make adjacent
rows agree.  Thus this encoding is an equivalence of existence predicates.

\subsection{The bounded-row CRT contract}

Let $E,F$ be polynomials with nonnegative integer coefficients in the row
index, fixed outer parameters, and $r$ local coordinates.  In the present
source $r=9$ and the coordinate bounds are
\[
 \mathbf u=(m,m,m,m,m,y,y,x,x).
\]
Suppose $t\ge1$ and
\[
 H\ge E(t,\mathbf u)+F(t,\mathbf u)+t+x+y+m+1.
\]
With the outer parameters suppressed, put
$T=H!$, $d_i=1+(i+1)T$, and
$\mathcal B=\prod_{i<t}d_i$.  Then the bounded family
\[
 \forall i<t\ \exists w\le\mathbf u:\ E(i,w)=F(i,w)
\]
is equivalent to the existence of natural $C,A_1,\ldots,A_r$ such that
\begin{align*}
 \mathcal B&\mid\binom{A_j}{u_j+1} &&(1\le j\le r),\\
 \mathcal B&\mid 1+(C+1)T, &
 E(C,A)&\equiv F(C,A)\pmod{\mathcal B}.
\end{align*}
The representatives $A_j$ may all be required positive, as they are in the
literal compiler.

For soundness, choose a prime $\ell\mid d_i$.  Since $d_i$ is coprime to
$H!$, necessarily $\ell>H$.  Multiplying the range congruence by
$(u_j+1)!$ gives a zero falling factorial modulo $\ell$, so
$A_j\equiv w_j\pmod\ell$ for some $0\le w_j\le u_j$.
The index congruence gives $C\equiv i\pmod\ell$ by cancellation of $T$.
The two evaluations are congruent and both lie in $[0,H]$, hence are equal.
For completeness, the $d_i$ are pairwise coprime: a common prime would
divide $j-i$ and so divide $H!$, which is impossible.  CRT encodes the
chosen row witnesses.  Each $(u_j+1)!$ divides $H!$ and is coprime to
$\mathcal B$, so it cancels from the falling-factorial congruence.
Positive representatives follow by adding multiples of $\mathcal B$.
The literal evaluation congruence uses
$E+\mathcal B q=F+\mathcal B v$ with natural $q,v$; arbitrary signed
quotients are represented by choosing the common offset $v$ sufficiently
large, so this orientation adds no restriction.

The supplied majorant certificate checks 392 gate inequalities and bounds
$E+F+t+x+y+m+1$ by $R_0^{6,893,174}$, where
$R_0=2+t+x+y+m+M_0\ge2$.  The literal choice
$H=R_0^{2^{23}}$ is therefore sufficient.  Each addition uses the exponent
bound $\max(e_1,e_2)+1$, each multiplication uses $e_1+e_2$; induction on
the actual nonnegative pair circuit proves the bound.  The factorial
bootstrap supplies the genuine value of $T$ before this CRT theorem is
invoked.

The fixed-length modulus product is itself represented arithmetically.
For $Q=(1+(t+1)T)^t+2$ and $D=TQ-1>0$,
\[
 \mathcal B=
 \left(T^t\,t!\binom{Q+t}{t}\right)\bmod D.
\]
Indeed the numerator is $\prod_{i=1}^t(TQ+iT)$, congruent to
$\prod_{i=1}^t(1+iT)$ modulo $D$.  The latter product is strictly smaller
than $D$ by the defining bound on $Q$.  An ordinary quotient/remainder
equation with remainder slack supplies this graph; no unbounded product
operator remains in the polynomial.

\subsection{One radix for all binomial requests}

For the twelve actual requests $\binom{n_j}{k_j}$, all $n_j,k_j$ are
positive.  Let $N$ be the sum of the distinct upper-argument circuit atoms
and $q=2^N+1$.  Thus $q>2^{n_j}$.  With
$L_j=q^{k_j}$ and $U_j=(q+1)^{n_j}$, the literal natural equations are
\[
 U_j=L_j(qh_j+C_j)+r_j,\qquad
 L_j=r_j+s_j+1,\qquad q=C_j+d_j+1.
\]
They recover exactly digit $k_j$ of the carry-free base-$q$ expansion of
$(q+1)^{n_j}$; therefore $C_j=\binom{n_j}{k_j}$, including $k_j>n_j$.
Conversely ordinary division supplies all four natural auxiliary values.
This theorem is simultaneous for arbitrary assigned argument values.
In completeness, the shared radix is selected after all semantic upper
arguments, including the chosen CRT representatives, have been fixed.

\subsection{Shared Pell certificates for all positive powers}

Write
$(a+\sqrt{a^2-1})^j=X_j(a)+Y_j(a)\sqrt{a^2-1}$ for $a\ge2$.
The imported number-theoretic ingredients are classification of natural
Pell solutions and the finite Matiyasevich index criterion.  The exact
repository interface is
\nolinkurl{DiophantineConsistency.PellRepresentation.criterion_iff_pell},
in \nolinkurl{provenance/upstream/calibration/PellRepresentation.lean},
which invokes \nolinkurl{Pell.matiyasevic} from Mathlib's
\nolinkurl{Mathlib/NumberTheory/PellMatiyasevic.lean}.

For the 23 distinct power requests $(n,k,m)$, the common frame is
\[
 W=2+\sum\{\hbox{distinct base, exponent, and output atoms}\},\quad
 Z=z+1,\quad A^2-W^3(W+2)Z^2=1.
\]
Each of the seven distinct bases has a natural $T_n$ satisfying
$2An=T_n+n^2+1$.  This equation excludes $n=0$ before any factorial or
power value is recovered.  Every exponent is syntactically positive.
Pell classification applied at base $W+1$ gives
$A=X_j(W+1)$ and $WZ=Y_j(W+1)$.
Modulo $W$, the recurrences give $Y_j(W+1)\equiv j$; thus
$j\ge W$.  Since $1\le n,k\le W<A$ and $m<W$, consequently
\[
 A\ge(W+1)^W>W^W,\qquad
 T_n\ge2A-2>W^W\ge n^k,\qquad m<T_n.
\]
Conversely $j=W$ supplies a frame for any finite family of intended
values.  The weaker inequality $A>W^2$ alone is insufficient for power
recovery; the divisibility of the frame's Pell index is essential.

For each of the eighteen exponent atoms $k_i$, use natural
$X_i,Y_i,S_i,R_i,e_i,s_i,r_i$.  The literal shared index equations are
\begin{align*}
 Y_i&=k_i+e_i,&X_i^2-(A^2-1)Y_i^2&=1,\\
 S_i^2-(B^2-1)R_i^2&=1,&S_i&=X_i+Us_i,\\
 R_i&=k_i+4Y_ir_i,&U^2-(A^2-1)V^2&=1,\\
 V&=\left(\prod_iY_i^2\right)(v_0+1),&
 B&=1+4V(b_0+1)=A+Ub_1.
\end{align*}
They imply every hypothesis of the nonzero branch of the imported index
criterion: $A,B>1$, $Y_i\ge k_i>0$, $V>0$, $Y_i^2\mid V$, and the four
required congruences.  Hence $(X_i,Y_i)=(X_{k_i}(A),Y_{k_i}(A))$.
For simultaneous completeness, set $L=4\prod_iY_i^2$ and choose a positive
power of the unit $A+s$ equal to one in the finite ring
$(\mathbb Z/L\mathbb Z)[s]/(s^2-(A^2-1))$.
Its Pell coordinates give $U\equiv1\pmod L$, $V\equiv0\pmod L$, $V>0$,
and $\gcd(U,4V)=1$.  CRT chooses $B>A$ with the displayed congruences.
Set $S_i=X_{k_i}(B)$, $R_i=Y_{k_i}(B)$; recurrence congruences and
monotonicity give all required natural quotients.

Finally each request has one natural quotient $h$ and
\[
 X_i+nY_i=AY_i+m+T_nh.
\]
For $H_j=X_j(A)-(A-n)Y_j(A)$, the recurrence is
$H_{j+1}=nH_j+T_nY_j(A)$, $H_0=1$.  Thus
\[
 H_k-n^k=T_n\sum_{j<k}n^{k-1-j}Y_j(A).
\]
Both $m$ and $n^k$ are in $[0,T_n)$, proving soundness; the displayed
nonnegative sum supplies $h$ for completeness.  The finite-ring and
recurrence arguments justify shared witnesses; merely identifying
independently existential variables would not suffice.

\paragraph{Composition inside ZFC.}
These proofs use finite induction, integer/rational arithmetic, prime
factorization, CRT, finite groups, the stated Pell theorems, and the
identified logical calibration.  They can be carried out over ZFC's own
$\omega$ and coded finite sequences, rather than only for externally
standard numerals.  After the factorial fixed-point proof and exact
polynomial projections are composed with these contracts, ZFC proves
existence of a polynomial zero equivalent to a coded ZFC refutation;
negating gives the asserted consistency biconditional.  Exact file checks
verify the supplied finite circuits and coefficients.  They do not prove
consistency or replace any imported semantic theorem.


\subsection{Order of soundness and completeness}

In soundness, first establish the natural domain and read the literal
arithmetic gates. The unconditional base-modulus rows exclude zero power
bases, and all exponents are syntactically positive. Apply the shared
power theorem, then the common-radix theorem. Only then use the
fixed-point theorem to identify $F_H$ and the shared-radius estimate to
identify $F_t$. The bounded CRT, beta, carrier, and flow arguments now
recover C222 and the calibrated refutation. In particular, the request
$F_H^t$ does not assume $F_H=H!$ in order to prove it.

In completeness, start from a finite proof. Choose its outer row and
label bounds, carrier, accumulator path, and beta data; these determine
$H$. Set the actual factorials $F_H=H!$ and $F_t=t!$ and their shared
radius. Next choose the CRT column and index representatives, then the
common binomial radix, then the true values of the finitely many powers.
Finally choose the common Pell frame and index auxiliaries. The polynomial
equations contain mutual references, but this construction of witnesses
has no unresolved cycle.

\section{From natural zeros to integer zeros with fewer guards}
\label{sec:integer}

The square-radius parent has 781 natural coordinates and 677 quadratic
residuals. Guarding each coordinate independently with four squares
would be wasteful: many signs are forced by rows already present.
The construction instead supplies a finite, ordered sign derivation for
every original coordinate. There is no inference from a variable's name
or intended interpretation.

\subsection{An exact ternary guard}

Define, for unrestricted integers $a,b,c$,
\[
                     \mathcal H(a,b,c)=a^2+b^2+c(c+1).
\]
This is nonnegative because $c(c+1)\ge0$ for every integer $c$.
Conversely, every $n\in\N$ is represented. Legendre's three-square
theorem applies to $4n+1$, which is not of the excluded form
$4^r(8s+7)$. Modulo four, exactly one of the three square roots is odd.
Write the roots as $2a,2b,2c+1$ and divide by four to obtain
$n=a^2+b^2+c(c+1)$. This representation is also explicitly formalized
in the AFP entry \emph{Three Squares}~\cite{three-squares}.

Appending $x=\mathcal H(a,b,c)$ is therefore an exact existential
nonnegativity guard. It adds three integer unknowns and one quadratic
row; it does not substitute a sum of squares into every occurrence of
$x$ and increase their degrees.

\subsection{Four sign rules}

\paragraph{Nonnegative polynomial definitions.}
Suppose a literal row says $x=f$, where $f$ has nonnegative integer
coefficients. Each monomial is nonnegative if every variable occurring
to an odd exponent in it has already been shown nonnegative. Even
powers need no sign assumption. The checker compares the complete
coefficient dictionary of the defining row and checks precisely these
dependencies.

\begin{lemma}[One guard for two coordinates]\label{lem:pair}
For $x,y\in\Z$ with $x+y\ge-1$,
\[
 x\ge0\text{ and }y\ge0
       \quad\longleftrightarrow\quad 2xy+x+y\ge0.
\]
\end{lemma}
\begin{proof}
The forward implication is immediate. Conversely,
$(2x+1)(2y+1)=2(2xy+x+y)+1>0$. The two odd factors have the same
sign. If both were negative, then $x,y\le-1$, contradicting
$x+y\ge-1$. Thus both are positive and $x,y\ge0$.
\end{proof}

When a literal sum row proves $x+y\ge-1$, append the single row
$2xy+x+y=\mathcal H(a,b,c)$. This is still quadratic and proves both
signs. It is not an attempt to make two independent square
representations share witnesses; it represents one nonnegative
centered product.

\begin{lemma}[A shifted quotient]\label{lem:quotient}
If $A\ge0$, $B\ge1$, $q\in\Z$, and $A(q+1)=B$, then $q\ge0$.
\end{lemma}
\begin{proof}
If $q\le-1$, the left side is nonpositive, a contradiction.
\end{proof}

For a binomial digit certificate, the literal rows unfold to
\[
 N=d(Qh+o)+r,\qquad d=r+s+1,\qquad Q=o+t+1.
\]
They give
\begin{equation}\label{eq:quotientidentity}
                     dQ(h+1)=N+dt+s+1.
\end{equation}
With all displayed dependencies nonnegative, Lemma~\ref{lem:quotient}
proves $h\ge0$ without a guard. For residuals
$R_1=N-d(Qh+o)-r$, $R_2=d-r-s-1$, and $R_3=Q-o-t-1$, the
difference of the two sides in \eqref{eq:quotientidentity} is exactly
$-R_1+dR_3+R_2$. Euclidean division has the analogous identity
$d(q+1)=N+s+1$. The checker expands and compares these finite
identities, including the necessary arithmetic definitions.

\begin{lemma}[Orient a root before Pell classification]\label{lem:frame}
Suppose the actual common frame gives
\[
 W=2+\sum\text{nonnegative atoms},\qquad
 A^2=1+W^3(W+2)(z+1)^2,
\]
with $A,z\ge0$. Then $A>W^2$. If $n\le W$ and
\[
                       X+Yn=YA+m+Tq
\]
with $Y,m,T,q\ge0$, then $X\ge0$.
\end{lemma}
\begin{proof}
Here $W\ge2$ and $(z+1)^2\ge1$, so
$A^2\ge1+W^3(W+2)>W^4$. As $A\ge0$, $A>W^2\ge n$.
Therefore $X=Y(A-n)+m+Tq\ge0$.
\end{proof}

This argument establishes the sign needed to apply the later Pell
theorem; it does not assume the Pell classification while trying to
prove one of its hypotheses. The certificate checks the exact frame
norm, the sum defining $W$, the relevant base-modulus row, and a
nonnegative-coefficient identity for $W-n$. All required signs occur
earlier in its schedule.

\subsection{The actual sign certificate}

The frozen construction has 163 ordinary guards and eight paired guards.
The latter seed 16 coordinate signs with only eight ternary
representations. It then uses the rules in Table~\ref{tab:signs}.
All 781 original coordinates are proved nonnegative. Guard auxiliary
coordinates remain unrestricted, as they should.

\begin{table}[htbp]
\centering
\begin{tabular}{@{}lr@{}}
\toprule
Rule applications after ordinary guard seeds & Number\\
\midrule
Nonnegative polynomial evaluation & 571\\
Shifted quotient & 13\\
Centered pair (two conclusions each) & 8\\
Common-frame root orientation & 18\\
\bottomrule
\end{tabular}
\caption{Exact sign-certificate census. Counting conclusions gives
$163+571+13+2\cdot8+18=781$.}
\label{tab:signs}
\end{table}

There are $171=163+8$ ternary guard rows, giving
$781+3\cdot171=1294$ unrestricted variables and
$677+171=848$ residuals before projection.

Starting from an arbitrary integer zero, the guards and the ordered
sign proof show that its original 781 coordinates form a natural zero
of the unchanged parent. Conversely, a natural parent zero makes
every selected coordinate and centered product nonnegative.
The ternary representation extends it to an integer zero of the guarded
system. This proves both existential directions without searching for
solutions or assuming that every sign is syntactically obvious.

\subsection{Certified monic projection}

Over $\Z$, a row $v-f(y)=0$ can be removed by substituting $f$ for $v$
in every other row. The inverse uniquely restores $v=f(y)$, regardless
of the signs of the coefficients of $f$. Each step checks that the
chosen row is exactly monic, that $f$ does not contain $v$, and that
\emph{every} remaining residual stays quadratic. Guards are substituted
along with all other rows and are never discarded merely because
their variable has been eliminated.

For $P_{956}$, 338 coordinates and rows disappear: 223 affine
restorations and 115 quadratic restorations, comprising 271 arithmetic
gates and 67 other witnesses. Thus
\[
              1294-338=956,\qquad848-338=510.
\]
For $P_{989}$, restricting the support of a candidate restoration to at
most four terms leaves 305 valid eliminations and therefore 989 variables
and 543 rows. The support threshold controls the size of the resulting
expression, not the mathematical validity of the projection.

The final sum of squares is expanded with exact integer arithmetic.
The collected dictionary contains degree-four terms and every counted
coordinate occurs. This proves the advertised exact degree, support,
and height. Composing the preceding two-way transformations with the
direct source calibration proves Theorem~\ref{thm:main}.

% TeX-ready mathematical fragment. Requires amsmath, amssymb, and booktabs.
\section{A short local majorant for nine positive variables}

A separate reduction gives two explicit nine-positive-variable equations:
\[
 (442,\,1\,520\,186),\qquad (456,\,1\,376\,378),
\]
where the entries are live binary arithmetic operations and exact total degree.
The second improves both measures over the reviewed
$475$-operation, degree-$1\,379\,730$ equation with the same original31 tuple.
The first is a tradeoff against the $433$-operation,
degree-$1\,628\,768$ compact54-sharp equation.  Neither changes the
nine-variable bound.  The new construction is a short local Laurent
majorant combined with a permissible successor coordinate.  No global
minimality or historical priority is asserted.

\subsection{The local lemma}

Let $A\ge2$, $C\ge2$, $i\ge1$ be integers and put
\begin{align*}
 J&=4A-5,& \delta&=A^2-1,&D&=1+\delta C^2,\\
 e&=iJC^2,&F&=1+\delta De^2,&
 T&=iC^2\bigl[AJ(2D-C^2)+3C^2-4\bigr].
\end{align*}
Then
\begin{equation}\label{eq:short-local-majorant}
                  T>2\sqrt{DF}>0.
\end{equation}
To prove this, set $t=C^2\ge4$. Exact expansion gives
\[
 T^2-4DF=i^2t^2Q(A,t)-4D,
\]
where
\begin{align*}
 Q(A,t)={}&20t(3t-4)A^3-t(83t-116)A^2\\
          &-10(7t^2-14t+8)A+109t^2-224t+116.
\end{align*}
Writing $A=x+2$, $t=y+4$, the coefficients of
$Q(x+2,y+4)-(x+2)^2$ are
\[
\begin{array}{c|rrr}
 &y^0&y^1&y^2\\ \hline
 x^0&1344&816&117\\
 x^1&3580&2188&318\\
 x^2&2975&1852&277\\
 x^3&640&400&60
\end{array}
\]
and are all positive. Consequently $Q(A,t)>A^2$ and
\[
 T^2-4DF>t^2A^2-4\bigl[(A^2-1)t+1\bigr]
       =A^2t(t-4)+4t-4\ge12.
\]
The defining expression gives $T>0$, since
$2D-t=2+(2A^2-3)t>0$, $J>0$, and $3t-4>0$.
This proves~\eqref{eq:short-local-majorant}.

The boundary condition is essential: at $(A,C,i)=(3,1,1)$ one obtains
$T^2-4DF=-308$. It is therefore supplied before the final equation is
assumed, by replacing its carrier input with $C=c+1$ for a positive
unknown $c$. Canonical Pell witnesses already have $C>1$, so this
restriction preserves existential completeness.

\subsection{The complete polynomial}

The nine positive unknowns are $b,g,h,s,w,c,i,j,n$. The two inherited
source prefixes define $R,N,\operatorname{cap}$ from their complete
literal coefficient tables. The delivered package includes both tables,
all four-port numerals, all fixed circuit leaves, and the exact prefix
maps. The whole compact54-sharp and original31 tuples are kept separate.
Both source prefixes prove $R,N\ge8$ and $b\le R,N$ on the ambient
positive domain.

Set $C=c+1$ and retain
\begin{align*}
 U&=N^2w,&Y&=N^2s,&M&=RY,&P&=2M^2U,\\
 K&=R+1+h(P-1),&A&=M(U+1),&B_1&=2R+1,&W_1&=bw,\\
 J&=4A-5,&Q&=3W_1C-2(W_1^2-1),\\
 a_0&=1+(P^2-1)K^2,&V_0&=PK,\\
 m&=K^2(4\operatorname{cap}-g)
       -16\operatorname{cap}(C-KY)^2,\\
 \Omega&=Jm,&E_0&=J(n+W_1^2)+Q.
\end{align*}
For $D,F,e$ as above define
\begin{align*}
 \Lambda&=CD+1-A,&G&=A+F\Lambda,\\
 Z&=B_1+(j-1)C,&H&=C+FZ,&I&=1+(G^2-1)H^2,\\
 V_1&=TGH,&a_1&=4DFI,&\Theta&=E_0-\Omega(V_0+V_1).
\end{align*}
The equation is the ordinary integer polynomial
\begin{equation}\label{eq:local-majorant-nine}
 \bigl[\Theta^2+\Omega^2(a_0-a_1)\bigr]^2
       -4\Omega^2a_0\Theta^2=0.
\end{equation}
All named quantities are polynomial abbreviations, and add no unknowns.
The circuit counts charge the successor $C=c+1$, scalar products, and
all squares. There are no divisions or variable exponents.

\subsection{The unbounded equivalence proof}

The source bounds give $A\ge33280$. Also $D,F,G,H>0$, $G>A$, and
$H>C\ge2$. Therefore $I<(GH)^2$, and the new lemma gives
$V_1>\sqrt{a_1}$. Independently,
$V_0^2-a_0=K^2-1>0$, so $V_0>\sqrt{a_0}$.
These inequalities hold before a zero is assumed.

At a zero, $E_0=Jn+(J-2)W_1^2+3W_1C+2>0$ excludes $\Omega=0$.
Factoring the real two-root norm yields signed roots
$r_0^2=a_0$, $r_1^2=a_1$ with
\[
 E_0=\Omega(V_0+V_1+r_0+r_1).
\]
Both strict majorants make the bracket positive for all sign choices.
Hence $m>0$. Only now do the inherited margin estimates give
$|C-KY|<K/2$, $C>K$, and $C>B_1+W_1$.
Moreover $A^2>P$, $D>C^2$, $F\ge A^2$, and $I>A^2$, whence
\[
 DFI>C^2A^4>K^2P^2>a_0.
\]
Thus $x=r_0+r_1=\Theta/\Omega$ is nonzero and rational. The isolation
identity $2xr_0=x^2+a_0-a_1$ makes both roots rational, hence integral.
The rational number $r_1/2$ has integer square $DFI$, so it too is
integral. The norm relation now gives $J\mid Q$.

The retained saturated Pell theorem recovers $C=\psi_A(B_1)$ from the
square $DFI$; all its range hypotheses have now been established.
Together with $a_0$ square, $m>0$, and $J\mid Q$, these are the exact
inherited ratio conditions, which recover the encoded source zero.

Conversely, a source zero supplies canonical ratio witnesses and compatible
saturated-Pell auxiliaries with $W_1^2+Q/J<2D$. Its carrier satisfies
$C=\psi_A(B_1)>1$, so $c=C-1$ is positive. Choose positive roots
$u_0=\sqrt{a_0}$ and $u_1=2\sqrt{DFI}$, and define
\[
 n=m(V_0+V_1+u_0+u_1)-W_1^2-Q/J.
\]
This is an integer. Since $F>D$, the lemma gives $V_1>2D$, so $n>0$.
Then $\Theta=\Omega(u_0+u_1)$ and
\eqref{eq:local-majorant-nine} holds. This proves both existential
directions. Negation and the parents' retained source/refutation
calibration give the Con(ZFC) biconditional, with neither consistency nor
inconsistency assumed. No new Lean, Rocq, or object-ZFC proof term is claimed.

\subsection{Cancellation-resolved exact degree}

Let $r=\deg R$, $\nu=\deg N$, $a=r+4\nu+2$, and $W=116$.
The unchanged source has an isolated positive leading kernel coefficient.
The shared degrees are
\begin{align*}
 \deg\Omega&=5a+2W+3,&\deg V_0&=3a+r+1,\\
 \deg E_0&=a+4,&\deg\delta_0&=3a+r+2,
\end{align*}
where $\delta_0=V_0^2-a_0$.
In weights $(a,1,1,r)$ on $(A,C,i,Z)$, the unique leading monomials are
\begin{align*}
 [T]_{\rm top}&=8A^4C^4i,\\
 [T^2-4DF]_{\rm top}&=60A^3C^8i^2,\\
 [GH]_{\rm top}&=256A^{14}C^{15}i^4Z,\\
 [V_1]_{\rm top}&=2048A^{18}C^{19}i^5Z,\\
 [\delta_1]_{\rm top}&=3932160A^{31}C^{38}i^{10}Z^2,
\end{align*}
where $\delta_1=V_1^2-a_1$. The last line follows from
\[
 \delta_1=(T^2-4DF)(GH)^2+4DF(H^2-1).
\]
Thus $\deg V_1=18a+r+24$ and $\deg\delta_1=31a+2r+48$.
Set
\[
 B_* = E_0^2-2E_0\Omega(V_0+V_1)+\Omega^2(\delta_0+\delta_1).
\]
The exact outer identity is
\[
 F_9=B_*^2+4\Omega^2V_0V_1B_*
       +4\Omega^4(\delta_0V_1^2+a_0\delta_1).
\]
Its four summand degrees are
\[
 82a+4r+8W+108,\quad72a+4r+8W+85,\quad
 59a+3r+8W+62,\quad57a+4r+8W+62.
\]
The first is strictly greatest and has positive nonzero leading coefficient.
The exact degree is therefore
\[
                        86r+328\nu+8W+272.
\]
For $(r,\nu)=(8799,2324)$ this is $1\,520\,186$;
for $(r,\nu)=(7943,2110)$ it is $1\,376\,378$.
The independent checker also computes nonzero specialized leading
coefficients modulo three primes using the actual fixed numerals.
Those coefficient witnesses supplement the universal identity and
unique-leader argument.

\subsection{Validation and comparison}

The independent standard-library verifier imports no builder. It verifies
all fixed numerals and their recipes, complete prefix maps, every new
register, local positivity and norm identities, topology, liveness, exact
operation counts, degree attainment, and full deterministic regeneration.
Normal and optimized executions agree. It rejects the invalid $C=1$
boundary and an altered majorant gate.

The operation breakdown is
\[
\begin{array}{c|rrrr}
 &+&-&\times&\text{total}\\ \hline
 \text{compact2082}&159&22&261&442\\
 \text{H1868}&155&22&279&456
\end{array}
\]
with exactly nine positive variables in each case. The whole coefficient
bindings and inherited ordinary proofs remain explicit. Finite circuit
checks are not substituted for the unbounded native/compiler/refutation
calibration, and no global simplicity optimum is claimed.


\section{Shorter exact radical norms}
\label{sec:norms}

The final contribution changes only an evaluator, not its polynomial or
zero set. It consequently transfers the parent's exact degree and
semantic equivalence without a new witness transformation. All actual
fixed numeral payloads are retained.

Define the formal radical norm
\[
 F_m(x;a_1,\ldots,a_m)
  =\prod_{\varepsilon\in\{\pm1\}^m}
      \left(x+\sum_{j=1}^m\varepsilon_j\sqrt{a_j}\right).
\]
Conjugation shows that this is an integer polynomial in $x,a_1,\ldots,a_m$.
To handle one rational radicand, use the cleared polynomial
\[
 \delta^{2^{m-1}}
 F_m(x;a_1,\ldots,a_{m-1},C/\delta).
\]
The following are identities of integer polynomial rings, including at
$\delta=0$ after clearing. Their literal circuits have neither radicals
nor division. Any nonzero-denominator hypothesis in the parent's
semantic interpretation is unchanged.

\subsection{A polarized three-root norm}

Set
\[
 z=x^2,\qquad G=\delta(z+a-b)-C,\qquad
 U=\delta^2az,\qquad V=\delta bC.
\]
Then
\begin{equation}\label{eq:norm3}
 \delta^4 F_3(x;a,b,C/\delta)
                  =[G^2+4(U-V)]^2-16UG^2.
\end{equation}
One can verify this by taking the two conjugates of the two-root norm,
or by comparison with the symmetric form
\[
 [J^2-4(T+R(a+b))]^2-64TRx^2,
\]
where $J=\delta(x^2-a-b)-C$, $T=\delta^2ab$, and $R=\delta C$.
The exact checker compares all 35 coefficients in the free variables
$x,a,b,C,\delta$. Sharing the displayed products gives 18 operations,
one fewer than the retained 19-operation implementation.

\subsection{A four-root norm with shared inner products}

Let $s=a+b$, $x_2=x^2$, and define
\begin{align*}
 q_0&=\delta(x_2+c-s)-C,&q_1&=-2\delta x,\\
 T&=\delta^2ab,&R&=\delta C,&z&=64TR,\\
 A_0&=q_0^2+cq_1^2-4(T+Rs),&A_1&=2q_0q_1,\\
 f_0&=A_0^2+cA_1^2-z(x_2+c),&f_1&=2(A_0A_1+zx).
\end{align*}
In the quadratic algebra $u^2=c$, the inner cleared three-root norm
at $x-u$ is $f_0+f_1u$. Its norm is therefore
\[
                              f_0^2-cf_1^2.
\]
This uses 38 operations rather than 43. All 495 free polynomial
coefficients are compared. The archive supplies the generic identity;
it does not claim an additional complete four-root endpoint in the
comparison table.

\subsection{Pair diagonal conjugates in the five-root norm}

Work in the biquadratic algebra $u^2=c$, $v^2=d$. If the inner
cleared norm is written
\[
                    f=f_0+2f_1u+2f_2v+2f_3uv,
\]
pair simultaneous sign reversal first. Direct multiplication gives
\[
 f(u,v)f(-u,-v)=G+4uvH,
\]
where
\[
 G=f_0^2+4(cd f_3^2-cf_1^2-df_2^2),\qquad
 H=f_0f_3-2f_1f_2.
\]
The other two conjugates give $G-4uvH$. Thus
\begin{equation}\label{eq:norm5core}
                  \operatorname{Norm}(f)=G^2-16cdH^2.
\end{equation}
Keeping the odd components divided by two is the useful computational
choice: there is no need to materialize all three factors of two before
taking the norm.

For completeness, the inner components are supplied by the following
finite polynomial formulas. Set $t=x+u+v$, $s=c+d$, and
\begin{align*}
 q_0&=\delta(x^2+s-a-b)-C,&q_3&=2\delta,&q_1&=q_3x,\\
 T&=\delta^2ab,&R&=\delta C,&z&=64TR,\\
 A_0&=q_0^2+s q_1^2+cd q_3^2-4(T+R(a+b)),\\
 A_1&=2q_1(q_0+dq_3),&A_2&=2q_1(q_0+cq_3),\\
 A_3&=2(q_0q_3+q_1^2).
\end{align*}
Then
\begin{align*}
 f_0&=A_0^2+cA_1^2+dA_2^2+cd A_3^2-z(x^2+s),\\
 f_1&=A_0A_1+dA_2A_3-zx,\\
 f_2&=A_0A_2+cA_1A_3-zx,\\
 f_3&=A_0A_3+A_1A_2-z.
\end{align*}
Indeed $q=\delta(t^2-a-b)-C=q_0+q_1u+q_1v+q_3uv$, and these
are exactly the components of the cleared three-root polynomial
$[q^2-4(T+R(a+b))]^2-zt^2$.

The resulting five-root circuit uses 84 operations rather than 93.
The checker compares all 20,349 coefficients of the free norm identity
and verifies exact substitution into the seven complete boundary
registers of each 13-integer parent. This proves equality of the
\emph{whole} parent and child polynomials, not agreement on sampled
inputs. It gives the four 13-variable improvements in
Table~\ref{tab:small}; \eqref{eq:norm3} gives the 11-positive-variable
improvement.

\section{Literal specification, provenance, and proof status}
\label{sec:artifacts}

\subsection{What makes these equations explicit}

The direct and integer equations are supplied twice: as literal
quadratic residual arrays and as fully collected quartic coefficient
dictionaries. A term is encoded as
\begin{lstlisting}
[coefficient, [zero_based_variable_indices]]
\end{lstlisting}
Repeated indices denote powers; an empty index list denotes a constant.
For example, $-3X_2^2X_7$ is \texttt{[-3,[2,2,7]]}. Every coefficient
is a literal signed integer. Variable names are annotations and do not
alter the polynomial. The archive's manifest authenticates every byte.

The small-variable equations are literal finite straight-line programs
with addition, subtraction, and multiplication. Every fixed leaf is
bound to the included exact numeral bytes, separately for each
coherent source tuple. They are not parameterized families awaiting
a choice of coefficients. Radical, factorial, power, finite-product,
and proof-calculus notation occurs only in explanations and proofs.

Printing 10,988 monomials in the article would obscure the definition
without making it more exact. The coefficient file, its residual
definition, the deterministic generator, and the independent expansion
checker give four mutually checkable specifications. The full file is
part of this submission, not an external dependency or an unspecified
computable constant.

\subsection{Important entry points in the archive}

\begin{longtable}{@{}p{0.43\textwidth}p{0.53\textwidth}@{}}
\toprule
Path & Content\\
\midrule\endhead
\file{README.md} & Results, exact domains, commands, and proof scope.\\
\file{report/main.tex} and \file{report/main.pdf} & This article, with its TeX inputs.\\
\file{direct/equations/parent781/} & Frozen square-radius natural parent, macro ledger, and bound certificates.\\
\file{direct/equations/natural467/} & Complete output-linear natural quartic and restoring maps.\\
\file{direct/equations/natural468/} & Complete square-radius natural quartic.\\
\file{direct/equations/natural469/} & Smaller-support natural alternative.\\
\file{integer/equations/integer956/} & Main integer quartic, guarded parent, sign proof, restoring maps, and expansion.\\
\file{integer/equations/integer989/} & Integer alternative with 26-bit coefficient height.\\
\file{nine/equations/} & Two complete nine-positive-variable polynomial DAGs.\\
\file{norms/equations/} & Five complete shorter polynomial evaluators.\\
\file{direct/provenance/upstream/} & Identified inherited direct arithmetic and source-calibration materials.\\
\file{BASELINES.md} & Exact comparison sources at the pinned repository snapshot.\\
\file{verify.py} & Portable release checks, including all component checks.\\
\file{checks/} & Release receipts and independent review scopes.\\
\file{MANIFEST.json} & SHA-256 inventory of the delivered files.\\
\bottomrule
\end{longtable}

\subsection{What the exact checks establish}

The release checks rebuild both natural parents byte for byte, reproduce
the natural contractions, replay every restoring substitution, and
independently collect every squared residual. They authenticate source
pins, actual fixed numeral bytes, and constructor and macro bindings.
The integer checker imports neither its builder nor its projection
optimizer: it independently verifies sign dependencies, coefficient
identities, guard coverage, projections, and the final expansion.
It rejects a missing pair guard, a false frame margin, an altered
restoring polynomial, and an altered final coefficient.

The nine-positive-variable checker verifies the exact majorant defect,
its positivity certificate, all new register formulas, norm identities,
literal source-prefix bindings, liveness, operation counts, and the
cancellation-resolved exact degree. Its nonzero leading coefficients
are also authenticated for the actual numeral specialization.
The norm checker proves free polynomial identities by coefficient
comparison and authenticates their substitutions into each complete
parent.

Normal Python and optimized Python runs are supplied. The mathematical
checks use explicit exceptions, not removable \texttt{assert}
statements as their acceptance condition. The component executables
need only the Python standard library and local packaged files.
No verification command searches for a solution of a consistency
equation.

\subsection{What the checks do not establish}

Finite algebra checks do not replace the unbounded equivalence proofs.
The factorial fixed-point proof, source reparameterization, integer
sign rules, and nine-variable majorant argument are ordinary
mathematical proofs formalizable in ZFC. Their interpretation as
consistency equations also uses the identified inherited compiler,
finite-flow, and C222-to-ZF-to-ZFC calibration theorems.

This release does not claim a new complete Lean or Rocq proof term,
nor a fully emitted derivation in an encoded object-language ZFC proof
system. The retained Lean sources identify the calibration interfaces;
the release's standard-library checkers do not invoke a proof-assistant
kernel or rebuild the entire upstream formal development. A reader
wanting that stronger assurance must formalize the new transformations
and replay the identified upstream development.

The distinctions are material: a correct coefficient transcript, a
universal mathematical equivalence, a proof-assistant theorem, and an
encoded internal proof object are different deliverables. This archive
supplies the first two and identifies the inherited dependencies of
the second. It does not label an experimental receipt as a proof of
$\Con(\ZFC)$.

\subsection{Reproduction}

From the extracted archive root, run
\begin{lstlisting}
python3 verify.py
python3 -O verify.py
sh build_report.sh
\end{lstlisting}
The first two commands run the packaged checks in a temporary copy,
so the verified delivered files are not rewritten. The report build
requires a conventional LaTeX installation with the packages listed
in \file{report/main.tex}. Mathematical validation does not require
LaTeX or internet access.

For a particular integer equation, the independent entry point is
\begin{lstlisting}
python3 integer/verify_integer.py \
  --parent direct/equations/parent781/equation.json \
  --equations integer/equations/integer956
\end{lstlisting}
For reconstruction, \file{integer/build_integer.py} accepts the
literal natural parent and an output directory; the support caps
16 and 4 give the two delivered integer variants. The source compiler
accepts \texttt{--radius square} or \texttt{--radius linear}.

\section{Assessment of the contribution}

The most reusable result is Theorem~\ref{thm:fixedpoint}. A large
independent precision parameter is often introduced in Diophantine
approximations to an integer-valued function. Here a polynomial in
the proposed output supplies that precision, and a two-stage lower
bound proves uniqueness of the fixed point. The gain is not just
the replacement of one displayed formula by another: one actual
positive-power request, its distinct base-modulus machinery, and
associated coordinates disappear from the emitted polynomial.

The rule-branch reduction attacks a different source of size: two
labels were kept as local bounded coordinates even though the
selected inference rule already determined them. Eliminating them
before the bounded-quantifier compiler saves more than eliminating
two variables at the end. The enlarged carrier label bound is essential
to that argument and is included explicitly.

The integer reduction then uses consequences of the literal rows to
avoid independent sign representations. The shifted-quotient and
common-frame arguments are semantic facts about an integer zero
of those rows; exact identities and an acyclic dependency schedule
make those facts independently checkable. The resulting equation
remains quartic throughout.

The smaller-variable improvements pursue a separate goal. The
nine-variable construction obtains a stricter majorant with controlled
cost and exact degree, while the norm identities preserve entire
polynomials and reduce their evaluation cost. These are concrete
improvements over specified parents, with all coefficient choices
fixed. None establishes a globally simplest consistency equation or
certifies that the ideas have never appeared elsewhere.

\begin{thebibliography}{99}
\raggedright

\bibitem{proveit}
Vladimir Reshetnikov and contributors,
\emph{ProveIt}, repository snapshot
\texttt{\RepoPin}.
\url{https://github.com/VladimirReshetnikov/ProveIt/tree/8f8fac484b976ba0ab75fd173282e838ca5a3b85}.
Relevant directories: \file{Computability/HilbertTenthProblem},
\file{Computability/DiophantineConsistency/con-zfc-diophantine-forms},
and \file{docs/incoming}.

\bibitem{synchronized}
ProveIt incoming submission,
\emph{ConZFC Synchronized Certificates},
\file{ConZFC_Synchronized_Certificates_2026-10-24.zip}, at the snapshot
in~\cite{proveit}. The present direct parent imports its shared
positive-power frame, product-index sharing, and direct source lineage.

\bibitem{batched}
ProveIt incoming submission,
\emph{ConZFC batched certificates},
\file{con_zfc_batched_certificates_2026-10-23.zip}, at the snapshot
in~\cite{proveit}. Source of the 1,690-integer direct baseline,
the degree-1,072 13-integer parent, and the monic projection adapter.

\bibitem{normparents}
ProveIt incoming submission,
\emph{ConZFC 13 Integer Norms},
\file{ConZFC_13_Integer_Norms_2026-10-24.zip}, at the snapshot
in~\cite{proveit}. Exact retained parents and notices are included
under \file{norms/upstream/}.

\bibitem{three-squares}
Anna Danilkin and Lo\"ic Chevalier,
\emph{Three Squares}, Archive of Formal Proofs, 3 May 2023.
\url{https://isa-afp.org/entries/Three_Squares.html}.
Includes the integer nonnegativity representation
$n=x^2+y^2+z^2+z$.

\bibitem{jones}
James P. Jones, \emph{Universal Diophantine equation},
Journal of Symbolic Logic \textbf{47} (1982), 549--571.
Background for the separate universal-polynomial line; it is not an
ingredient of the new direct C222 quartic source.

\bibitem{universal}
Jonas Bayer, Marco David, Malte Ha\ss ler, Yuri Matiyasevich, and
Dierk Schleicher,
\emph{Diophantine equations over $\Z$: universal bounds and parallel
formalization}, arXiv:2506.20909, version 3, 22 October 2026.
\url{https://arxiv.org/abs/2506.20909}.
Background on keeping variable, degree, coefficient, and formalization
claims separate.

\end{thebibliography}
\end{document}
